% Irodov Part III §3.3 final problems 141--146;
% source PDF pp.120--121 (printed 111--112); figures 3.33, 3.34.
\electrocapproblem{141}{111}
空气平板电容器的极板面积均为 $S$，
缓慢把板间距从 $x_1$ 增至 $x_2>x_1$。
求所需的外界机械功：
(a) 极板自由电荷 $q$ 保持不变；
(b) 两极板电压 $V$ 保持不变。
\exsolution
忽略边缘场，
$C(x)=\varepsilon_0S/x$。
(a) 固定自由电荷时
$U(q,x)=q^2/(2C)=q^2x/(2\varepsilon_0S)$。
外界准静态功等于场能增加量：
\[
\boxed{A_q=\frac{q^2(x_2-x_1)}{2\varepsilon_0S}}.
\]
(b) 固定电压时电源会交换电荷，
不能简单取储能差作外功。
极板间吸引压力为
$p_e=\varepsilon_0V^2/(2x^2)$，
力 $F=p_eS$，
因此
\[
A_V=\int_{x_1}^{x_2}\frac{\varepsilon_0SV^2}{2x^2}
\,\mathrm dx
=\boxed{\frac{\varepsilon_0SV^2}{2}
\left(\frac1{x_1}-\frac1{x_2}\right)}.
\]
\excheck 原书(a)英文正文有
“capacitance ... equal to q”的明显笔误，
$q$ 量纲为电荷，依物理意义必须是
“charge $q$ remains constant”；
与原书答案的 $q^2$ 表达式一致。

\electrocapproblem{142}{111}
平板电容器原电容 $C_0=20\,\mathrm{nF}$，
极板间隙为 $d$。
先插入厚度为 $h=\eta d$ 的平行板
（$\eta=0.60$），接电压
$V=200\,\mathrm V$ 充电，随后切断电源，
再缓慢抽出该插板。
求抽出所需机械功：
(a) 金属板；(b) 玻璃介质板。
\exsolution
设充电时插板在位的电容为 $C_i$。
电源断开后自由电荷
$Q=C_iV$ 保持不变；
抽出后电容回到 $C_0$。
因此机械功
\[
A=\Delta U=
\frac{Q^2}{2}\left(\frac1{C_0}-\frac1{C_i}\right)
=\frac{C_i^2V^2}{2}
\left(\frac1{C_0}-\frac1{C_i}\right).
\]
(a) 金属板内部电场为零，
只剩有效气隙厚度
$d-h=d(1-\eta)$，
故 $C_i=C_0/(1-\eta)$，得到
\[
\boxed{A_{\rm metal}=
\frac12C_0V^2\frac{\eta}{(1-\eta)^2}
=1.5\,\mathrm{mJ}}.
\]
(b) 玻璃相对介电常数为 $\varepsilon$，
与两侧空气为法向串联。
设
$f=1-\eta+\eta/\varepsilon$，
则 $C_i=C_0/f$，所以
\[
\boxed{A_{\rm glass}=
\frac12C_0V^2
\frac{\eta(1-\varepsilon^{-1})}
{(1-\eta+\eta/\varepsilon)^2}}.
\]
\excheck 第(b)问题面未给玻璃 $\varepsilon$ 数值，
应保留参数表达式，
不可擅自代入某种玻璃的介电常数。

\electrocapproblem{143}{111}
将水平放置的平板电容器浸没在水中，
极板间距 $d=1.0\,\mathrm{mm}$，
水完全填满间隙。
两极板电压保持为 $V=500\,\mathrm V$。
求水在间隙中的电压力增量。
\exsolution
取水的相对介电常数为
$\varepsilon$，忽略边缘效应。
水中电场 $E=V/d$。
界面电场能和介质极化在固定外加电势的
电机械耦合关系给水中附加的电压力
\[
\Delta p=\frac12\varepsilon_0
\varepsilon(\varepsilon-1)E^2.
\]
代入 $E=V/d$：
\[
\boxed{\Delta p=
\frac{\varepsilon_0\varepsilon(\varepsilon-1)V^2}{2d^2}}.
\]
室温水取 $\varepsilon\simeq81$：
\[
\boxed{\Delta p\simeq7.2\times10^3\,\mathrm{Pa}
\approx7\,\mathrm{kPa}\approx0.07\,\mathrm{atm}}.
\]
\excheck 该式是本题所采用的水介电响应
电压与压力模型，
与只计算真空 Maxwell 电压力
$\varepsilon_0E^2/2$ 不同。

\electrocapproblem{144}{111}
水平平板电容器的一块极板浸在液体内，
另一块位于液面上方（图3.33）。
液体密度 $\rho$，相对介电常数
$\varepsilon$。
极板表面自由电荷密度模为 $\sigma$，
求液体在两极板之间升高的高度 $h$。
\exsolution
两板间浸入介质的部分增加时，
其单位体积静电储能由
$u_{\rm vac}=\sigma^2/(2\varepsilon_0)$
变为
$u_{\rm diel}=\sigma^2/(2\varepsilon_0\varepsilon)$。
在保持自由电荷密度的准静态近似下，
每单位上升体积可释放的电场能为
\[
\Delta u=\frac{\sigma^2}{2\varepsilon_0}
\left(1-\frac1\varepsilon\right).
\]
与上升液柱静水压头 $\rho gh$ 平衡，
\[
\rho gh=\Delta u.
\]
因此
\[
\boxed{h=\frac{(\varepsilon-1)\sigma^2}
{2\varepsilon_0\varepsilon\rho g}}.
\]
\excheck 如果改为固定\emph{极板电压}，
电荷密度会随介质分布改变，
必须使用相应不同的热力学约束。

\electrocapproblem{145}{111}
同轴圆柱电容器，
两电极平均半径为 $R$，
径向间隙 $d\ll R$。
向间隙沿轴向插入相对介电常数
$\varepsilon$ 的圆柱形介质层，
两电极保持恒定电压 $V$。
求把介质进一步吸入的电力。
\exsolution
忽略端部效应，插入轴向长度为 $x$
时增加的介质体积为
$2\pi Rd\,x$。
这相当于单位插入长度新增并联电容
\[
\frac{\mathrm dC}{\mathrm dx}
\simeq\frac{\varepsilon_0(\varepsilon-1)}d
(2\pi R).
\]
在固定电压下，
电源—电场系统的有效势为
$-\frac12CV^2$，
故吸入力
\[
F=\frac12V^2\frac{\mathrm dC}{\mathrm dx}.
\]
于是
\[
\boxed{F=\frac{\pi R\varepsilon_0
(\varepsilon-1)V^2}{d}}.
\]
\excheck 这是介质已部分插入时
接近轴向端缘的电力，
不表示介质无限延伸后的
横向均匀区域仍有非零吸入力。

\electrocapproblem{146}{112}
图3.34：固定电极由两片半径为 $R$
的半圆形金属极板组成，
板间距 $d$。薄可转动介质板绕
半圆中心 O 转动，
相对介电常数为 $\varepsilon$，
转角为 $\phi$，固定电压 $V$。
求绕 O 的电力矩模。
\exsolution
忽略边缘场，
介质覆盖的两极板重叠面积随转角改变。
半径 $R$ 的圆扇形面积为
$A(\phi)=\frac12R^2\phi$，
所以
$\mathrm dA/\mathrm d\phi=R^2/2$。
被介质覆盖区域相对于真空的
增量电容为
$\Delta C=(\varepsilon-1)\varepsilon_0A/d$，
因此
\[
\frac{\mathrm dC}{\mathrm d\phi}
=\frac{\varepsilon_0(\varepsilon-1)R^2}{2d}.
\]
固定电压的电力矩
$\mathcal N=\frac12V^2\mathrm dC/\mathrm d\phi$，
故
\[
\boxed{|\mathcal N|=
\frac{\varepsilon_0(\varepsilon-1)R^2V^2}{4d}}.
\]
\excheck 在原图允许的旋转角度范围内，
扇形覆盖面积随 $\phi$ 线性变化，
故电力矩大小不依赖 $\phi$。
